\chapter{Plan for the Next Semester}
\label{ca:chap:plan}

\section{Work completed this semester}

Table~\ref{ca:tab:done} summarises the work so far. Every run is launched from a job queue
that records its configuration, log and result, and every design decision and its evidence
is recorded in a decision log, so each number in this report can be traced to a run.

\begin{table}[htbp]
\centering
\small
\begin{tabular}{@{}l>{\raggedright\arraybackslash}p{8.6cm}r@{}}
\toprule
\textbf{When} & \textbf{Work} & \textbf{GPU-h} \\
\midrule
Sep            & Literature review; problem formulation and hypotheses & -- \\
to 4 Oct       & Motivating cross-tokenizer experiment (SoC cluster) & -- \\
10--11 Oct     & Server environment; provider screening (9 models, 3 domains); domains redesigned; design fixed & 3.0 \\
11 Oct         & Response cache (5 providers, 30,000 responses each) & 8.1 \\
13--18 Oct     & Pilot (11 runs); revised selectors; learning-rate check; simulator & 16.5 \\
17--24 Oct     & Main runs (30), checkpoint evaluations (9), branches (216), test evaluation & 80 \\
\bottomrule
\end{tabular}
\caption{Work completed, with L40 GPU-hours on the project server.}
\label{ca:tab:done}
\end{table}

\section{Next steps}
\label{ca:sec:next}

\begin{enumerate}[itemsep=3pt,topsep=3pt]
  \item \textbf{Settle H1 and H2.} Re-evaluate the 216 branch adapters, saving per-item
    results, so that providers can be compared on the same items within each seed. Then
    run longer branches (40 to 60 rounds) and more seeds on Countdown, sized by a power
    calculation from the current variances, for the providers whose values are close.
  \item \textbf{Complete H3.} Run the two further budgets, so that the comparison holds at
    more than one operating point. Add the planned selectors from the literature
    (cost-aware EXP3, sliding-window UCB and LinUCB), a static allocation that gives each
    domain its best provider, and an oracle approximation from the branch measurements.
  \item \textbf{Understand the training instability.} Diagnose the response-length
    collapse on Countdown, for example by tracking entropy and testing a KL penalty or a
    length floor, since one collapsed run decides the ranking of a selector.
  \item \textbf{Widen where learning happens.} Test whether a longer budget or a larger
    student produces learning on MATH and code, so that H1 can be tested across domains
    rather than mostly within Countdown.
  \item \textbf{Selection rewards and allocation.} Compare cheap proxies (loss decrease,
    learnability, signal density) with utility gain as the selector's feedback, and compare
    each provider's share of the budget with a post hoc valuation of its contribution.
  \item \textbf{External check against FedTOPS.} Evaluate the final students on the six
    mathematics benchmarks used by FedTOPS \autocite{chen_fedtops_2026}, as an
    out-of-distribution check that places the results next to the work that inspired the
    project.
  \item \textbf{Write the dissertation}, keeping the main text well within the 55-page
    limit and moving implementation detail to appendices.
\end{enumerate}

\begin{table}[htbp]
\centering
\small
\begin{tabular}{@{}lccccc@{}}
\toprule
\textbf{Task} & \textbf{Dec} & \textbf{Jan} & \textbf{Feb} & \textbf{Mar} & \textbf{Apr} \\
\midrule
1. Settle H1 and H2           & \textbullet & \textbullet &             &             &   \\
2. Complete H3                &             & \textbullet & \textbullet &             &   \\
3. Training instability       & \textbullet & \textbullet &             &             &   \\
4. Learning on MATH and code  &             & \textbullet & \textbullet &             &   \\
5. Selection rewards, allocation &          &             & \textbullet & \textbullet &   \\
6. FedTOPS benchmarks         &             &             & \textbullet &             &   \\
7. Dissertation and presentation & \textbullet &          &             & \textbullet & \textbullet \\
\bottomrule
\end{tabular}
\caption{Schedule for the next semester. December covers the vacation.}
\label{ca:tab:schedule}
\end{table}

\section{Resources and risks}
\label{ca:sec:resources}

\paragraph{Resources.} The project uses a shared server with six NVIDIA L40 GPUs of 48\,GB,
no job scheduler, and at most two GPUs at a time; the motivating experiment used the School
of Computing's Slurm cluster. About 110 L40 GPU-hours have been used. The plan above needs
an estimated 190 more: about 100 for two further budgets, 45 for the new selectors, 35 for
longer branches and branch re-evaluation, and 10 for the FedTOPS benchmarks. At two GPUs
this is roughly four days of continuous running, so compute is not the binding constraint.
Software: Python 3.12, PyTorch, vLLM, TRL and PEFT, with code run in a bubblewrap sandbox.

\paragraph{Risks.} The main risk is that differences between providers stay below what the
design can detect even with more seeds. That outcome would itself answer the research
question for this setting: price, not provider identity, decides value. The second risk is
that learning stays confined to Countdown, which would narrow the claims to one domain. The
third is GPU contention on the shared server, which the queue handles by waiting for idle
GPUs.
